\section{Objetivos}

\subsection{Objetivo general}

Sustentar un caso de uso de machine learning concurrente sobre un dataset real de más de un millón
de registros, fundamentado en literatura académica reciente. El caso servirá de base para
implementar en Go, sin librerías de terceros, una versión secuencial y otra concurrente de K-means
en los entregables siguientes.

\subsection{Objetivos específicos}

\begin{itemize}
  \item Revisar nueve artículos académicos publicados entre 2021 y 2026, tres por integrante, sobre
        modelos de machine learning paralelizados. En cada uno, identificar qué mecanismo de
        concurrencia usa y qué beneficio le aporta al modelo.
  \item Organizar la revisión en tres temas con preguntas aplicables a nuestra implementación: cómo
        repartir y reducir el trabajo en memoria compartida, cuánto cuesta sincronizar entre nodos y
        cuándo conviene hacer menos trabajo en lugar de repartirlo.
  \item Justificar la elección de K-means sobre los viajes de taxi de Nueva York y su relación con
        el ODS 11.
  \item Seleccionar y limpiar un mes completo del dataset con reglas explícitas, verificables y
        reproducibles, dejando más de un millón de registros válidos.
  \item Preparar las features y el formato de entrada (CSV) que consumirá la implementación en Go.
  \item Evidenciar el trabajo de todos los integrantes en GitHub con el flujo Git Flow.
\end{itemize}
